% TOP_PROBLEM: {"id":30007539,"problem_number":"LOCAL-30007539","title":"State-contingent sovereign borrowing","category_id":21,"category":{"id":21,"name":"economics","display_name":"Economics","description":"Open economic theory statements, published research agendas, and proposed scoped specifications; question provenance and historical priority are qualified per record.","slug":"economics","order_index":21,"created_at":"2026-10-08T00:00:00Z"},"status":"research_question","external_url":null,"legacy_ids":[],"published":false,"statement":"Determine when an indexed sovereign contract improves welfare after accounting for reporting incentives, lender participation, ambiguity, and liquidity costs.","economics":{"original_id":28,"field":"Fiscal policy and sovereign debt","kind":"design","formalization_basis":"adapted_research_question","source_status":"explicit_open","priority_status":"earliest_found","priority_note":"This is the earliest explicit relevant statement verified in this audit, not a claim of original priority; older literature and earlier versions may contain antecedents.","earliest_verified_reference":{"authors":["Leonardo Martinez","Francisco Roch","Francisco Roldán","Jeromin Zettelmeyer"],"title":"Sovereign Debt","year":2022,"url":"https://fqroldan.github.io/resources/MRRZ.pdf","doi":"10.5089/9798400213250.001","locator":"Section 10, “Policy options to reduce sovereign risk,” subsection “State-contingent debt,” pp. 25–26","evidence_summary":"The survey calls the lack of indexation a puzzle despite theoretical benefits, notes little theoretical treatment of implementation obstacles, and reviews competing explanations. It cites earlier work including Bolton and Jeanne (2009), whose original puzzle wording was not verified here."},"current_reference":{"authors":["Leonardo Martinez","Francisco Roch","Francisco Roldán","Jeromin Zettelmeyer"],"title":"Sovereign Debt","year":2022,"url":"https://fqroldan.github.io/resources/MRRZ.pdf","doi":"10.5089/9798400213250.001","locator":"Section 10, “Policy options to reduce sovereign risk,” subsection “State-contingent debt,” pp. 25–26","evidence_summary":"The survey calls the lack of indexation a puzzle despite theoretical benefits, notes little theoretical treatment of implementation obstacles, and reviews competing explanations. It cites earlier work including Bolton and Jeanne (2009), whose original puzzle wording was not verified here."},"formalization_note":"The source explicitly calls limited indexation a puzzle. This reporting/liquidity mechanism is a representative specification, not a claim that the survey first posed the puzzle."},"statusQualification":"Published question or research agenda verified at the cited locator. The mathematical specification is adapted for this catalogue and includes additional assumptions; source publication does not establish first-ever priority or certify this exact model remains unsolved. The source explicitly calls limited indexation a puzzle. This reporting/liquidity mechanism is a representative specification, not a claim that the survey first posed the puzzle.","statusReviewedAt":"2026-10-08"}
% FIRST_POSTED: 2026-10-08
% ENTRY_KIND: statement_only
% !TeX program = lualatex
\documentclass[11pt]{article}
\usepackage{fontspec}
\usepackage[a4paper,margin=25mm]{geometry}
\usepackage{amsmath,amssymb,mathtools}
\usepackage{xurl}
\usepackage[hidelinks]{hyperref}
\newcommand{\sref}[2]{\hyperlink{source:#1:#2}{[S#2]}}
\setlength{\emergencystretch}{3em}
\begin{document}
% problemId:30007539
\section*{State-contingent sovereign borrowing}\label{30007539}
\subsection{Definitions and mathematical statement}
Fix a sovereign with stochastic true output \(y_t\), observable output report \(\tilde y_t\), and an ability to influence reporting at a real cost. A debt contract promises repayment \(g(\tilde y_t)\), possibly subject to a maximum payment. Investors have a declared stochastic discount factor, ambiguity set for output and reporting, and a liquidity cost depending on contract standardization. The sovereign’s preferences and default technology are given. Let \(\mathcal G\) be a finite-dimensional, publicly stated family of repayment schedules.
Determine whether a schedule \(g\in\mathcal G\) yields a welfare-improving equilibrium relative to noncontingent debt while satisfying lender break-even and truthful-reporting incentives. In particular, require
\[
E[m_t g(\tilde y_t)]-\text{liquidity and ambiguity costs}\ge \text{issue price},
\]
where \(m_t\) is the investor discount factor, together with a specified sovereign incentive-compatibility condition for every feasible false report. Quantify how much insurance survives pricing premia, manipulation, and reduced market liquidity; explain conditions under which noncontingent issuance is an equilibrium despite gross risk-sharing benefits. Assess contract standardization or independent verification as feasible remedies. The problem concerns this explicitly modeled adoption puzzle, not whether state-contingent debt can provide insurance in an ideal frictionless model, a benefit already established.

\par\textbf{Formulation and scope.} The source explicitly calls limited indexation a puzzle. This reporting/liquidity mechanism is a representative specification, not a claim that the survey first posed the puzzle.

\par\textbf{Open-question provenance.} An explicit question or research agenda was verified at the source locator. This does not by itself certify that every later version remains unresolved \sref{30007539}{1}.

\par\textbf{Historical priority.} This is the earliest explicit relevant statement verified in this audit, not a claim of original priority; older literature and earlier versions may contain antecedents.

\subsection{Short English statement}
Determine when an indexed sovereign contract improves welfare after accounting for reporting incentives, lender participation, ambiguity, and liquidity costs.

\subsection{Sources}
\begin{itemize}\raggedright
\item[\textnormal{[S1]}]\hypertarget{source:30007539:1}{} Leonardo Martinez, Francisco Roch, Francisco Roldán, Jeromin Zettelmeyer. 2022. Sovereign Debt. Section 10, “Policy options to reduce sovereign risk,” subsection “State-contingent debt,” pp. 25–26.
\url{https://fqroldan.github.io/resources/MRRZ.pdf}.
DOI: 10.5089/9798400213250.001.
{\small\textit{Earliest verified question/agenda reference; historical priority qualified below. The survey calls the lack of indexation a puzzle despite theoretical benefits, notes little theoretical treatment of implementation obstacles, and reviews competing explanations. It cites earlier work including Bolton and Jeanne (2009), whose original puzzle wording was not verified here.}}
\end{itemize}
\end{document}
